% Layout source for the generated book. Educational prose: CC BY-SA 4.0; TeX layout code: MIT.
\documentclass[11pt,oneside,openany]{book}
\usepackage[T1]{fontenc}
\usepackage{lmodern}
\usepackage[letterpaper,inner=1.0in,outer=1.0in,top=0.9in,bottom=0.95in,headheight=15pt]{geometry}
\usepackage{amsmath,amssymb}
\usepackage{xcolor,array,longtable,enumitem}
\usepackage{fancyhdr}
\usepackage{hyperref}
\definecolor{AlgebraBlue}{HTML}{244A73}
\definecolor{TrickPurple}{HTML}{6F398E}
\definecolor{Ink}{HTML}{172D40}
\hypersetup{colorlinks=true,linkcolor=AlgebraBlue,urlcolor=AlgebraBlue,pdftitle={Calculus: A Worked Review},pdfauthor={}}
\setlength{\parindent}{0pt}
\setlength{\parskip}{6pt plus 2pt}
\setlength{\emergencystretch}{3em}
\setlist{itemsep=4pt,topsep=4pt,leftmargin=*}
\setcounter{secnumdepth}{0}
\setcounter{tocdepth}{1}
\pagestyle{fancy}
\fancyhf{}
\fancyhead[L]{\small\sffamily Calculus: A Worked Review}
\fancyhead[R]{\small\sffamily\thepage}
\fancyfoot[L]{\small\sffamily Alpha 0.2 | Notation and worked-example prototypes}
\renewcommand{\headrulewidth}{0.3pt}
\renewcommand{\footrulewidth}{0pt}
\begin{document}
\frontmatter
\begin{titlepage}
\color{Ink}
\vspace*{0.7in}
{\sffamily\large A PERSONAL PROBLEM-SOLVING REFERENCE\par}
\vspace{0.5in}
{\sffamily\bfseries\fontsize{40}{46}\selectfont Calculus\par}
\vspace{0.15in}
{\sffamily\LARGE A Worked Review\par}
\vspace{0.45in}
{\Large Explicit notation. Visible reasoning.\par}
{\Large Every dependency accounted for.\par}
\vfill
{\sffamily\large Alpha development edition 0.2\par}
\medskip
Notation standard, curriculum, eight worked problems, and supporting tools.\\
Website foundation prepared October 7, 2026.
\end{titlepage}
\chapter*{How this book will work}
This is a review book for readers who have encountered calculus but want a dependable way to approach unfamiliar problems. Every object has a declared type. Every function shows its inputs. Every derivative and integral identifies the variable being used.

Worked solutions expose the choices as well as the calculations. Algebra and identities carry explicit labels; dedicated algebra steps appear in \textcolor{AlgebraBlue}{dark blue}. Supporting tools from the tricks appendix appear in \textcolor{TrickPurple}{purple}. Color supplements the labels, so a grayscale copy remains usable.

This alpha development edition establishes the notation with a scalar differentiation chapter, two further pilot problems, and a supporting tricks appendix. It does not yet contain the full curriculum. The proposed main parts are Differentiation and differentials, and Integration. The notation recommendations incorporate a targeted literature review, documented in the references.

Try a question before reading its hints. Use the stable problem ID and step number when asking a tutor for help. The tutor should keep this notation and work from the last step you understand.
\chapter*{Sharing and accuracy}
Copyright \copyright{} 2026 Calculus Workbook contributors, to the extent copyright applies. Original educational content is licensed under Creative Commons Attribution-ShareAlike 4.0 International (CC BY-SA 4.0). You may share and adapt it, including commercially, subject to attribution, an indication of changes, ShareAlike, and the license's other terms.

\url{https://creativecommons.org/licenses/by-sa/4.0/}

This developing book is provided as-is. Mathematical correctness, completeness, and suitability for a particular purpose are not guaranteed. The license's warranty disclaimer and liability limits apply to the extent permitted by law. Corrections are welcome. Linked references retain their own rights.

\tableofcontents
\mainmatter

\chapter{The notation standard}

Status: recommended standard, version 0.2, revised October 5, 2026 after a targeted literature review. This is a consistent selection of established conventions, with explicit teaching annotations. It is not a claim that all fields use one universal notation.

\section{1. Reading contract}

Use \textbf{Leibniz operator notation with full arguments}, ordinary parentheses for grouping, and explicit type and dependency declarations. Avoid inventing symbols to replace established calculus notation. The literature basis and reasons for the choices are recorded in the literature review.

Every problem and solution begins with four items:

\begin{enumerate}
\item \textbf{Objects and types:} scalar, vector with dimension, matrix with dimensions, or tensor with order, space, and component shape.
\item \textbf{Dependencies:} independent inputs, dependent outputs, intermediate functions, and fixed parameters.
\item \textbf{Domain and assumptions:} allowed inputs, excluded values, and regularity or sign conditions used.
\item \textbf{Operation and output:} what is being differentiated or integrated, with respect to which variable, what stays fixed, and the output type.

\end{enumerate}
Repeat declarations in the solution so it can be read independently. Full function arguments stay visible in calculus expressions. A vector or matrix argument may be written as a single object only when its components and dependencies are displayed locally. Long expressions are split into named, fully defined steps rather than compressed into unexplained symbols.

\section{2. Delimiters, types, and dependencies}

\subsection{Give delimiters predictable jobs}

\begin{longtable}{@{}>{\raggedright\arraybackslash}p{0.300\textwidth}>{\raggedright\arraybackslash}p{0.300\textwidth}>{\raggedright\arraybackslash}p{0.300\textwidth}@{}}
Notation & Meaning in this book & Safeguard \\[5pt]
\hline\endhead
\(f(x,y)\) & Function value at the listed inputs & Parentheses immediately after a function name contain its arguments \\[5pt]
\((a+b)\) & Grouping & Larger parentheses group the entire operand of a derivative \\[5pt]
A displayed bracketed column or rectangular array & Vector or matrix representation & State its type and shape; do not use a one-line tuple as an unexplained vector \\[5pt]
\([a,b]\) & Closed interval & Preserve this standard meaning; in first uses also write \(a\leq x\leq b\) \\[5pt]
\(\{x: x>0\}\) & Set & Braces are not the default grouping delimiter \\[5pt]
\(A_{ij}(t)\), \(T_{ijk}(t)\) & Scalar components & Declare the parent object, basis, index ranges, and dependencies \\[5pt]
\end{longtable}

Square brackets do not universally mean “tensor.” They also conventionally denote intervals and sometimes operator arguments. Our deliberate local convention is to avoid square brackets for derivative operands and ordinary algebraic grouping. We retain them for arrays and closed intervals because changing those conventions would introduce a different confusion.

For nested grouping, use parentheses of different sizes. If a formula becomes hard to track, break it into displayed steps. Delimiters alone never determine an object's type.

\subsection{Object typography}

\begin{longtable}{@{}>{\raggedright\arraybackslash}p{0.300\textwidth}>{\raggedright\arraybackslash}p{0.300\textwidth}>{\raggedright\arraybackslash}p{0.300\textwidth}@{}}
Object & Convention & Example declaration \\[5pt]
\hline\endhead
Scalar & Ordinary italic & \(x\in\mathbb R\): independent real scalar \\[5pt]
Scalar-valued function & Ordinary letter, all inputs visible & \(z=f(x,y)\in\mathbb R\): dependent scalar output \\[5pt]
Vector & Bold lowercase & \(\mathbf r(t)\in\mathbb R^3\): column vector depending on scalar \(t\) \\[5pt]
Matrix & Bold uppercase & \(\mathbf A(t)\in\mathbb R^{m\times n}\): matrix depending on scalar \(t\) \\[5pt]
Tensor of order three or higher & Bold calligraphic uppercase & \(\boldsymbol{\mathcal T}(t)\): tensor with declared space, order, basis, and component shape \\[5pt]
\end{longtable}

These typography conventions are established choices, not universal rules. Scalar entries use ordinary type: \(r_i(t)\), \(A_{ij}(t)\), and \(T_{ijk}(t)\). Tensor order is not called tensor rank: rank has other meanings.

For this book, a function name and a function value are distinct: \(f\) names the function; \(f(x,y)\) supplies its inputs. We do not write \(z=f\) when we mean \(z=f(x,y)\).

Fixed parameters appear after a semicolon, as in \(f(x;a)=ax^2\), and are also described in words. This semicolon is an organizational convention, not a different calculus operation. If a parameter starts varying, announce that fact and write its dependence explicitly.

For independent scalar coordinates \(x,y\) and a dependent scalar \(z=f(x,y)\), say that \(x\) and \(y\) vary independently. For motion along a path, display instead

\[
t\longmapsto (x(t),y(t))\longmapsto z(t)=f(x(t),y(t)).
\]

Here \(t\) is the independent scalar parameter; \(x(t)\) and \(y(t)\) are dependent scalar coordinates; \(z(t)\) is a dependent scalar output. The parenthesized coordinate pair is a point, not an undeclared column vector.

\section{3. Derivatives}

\subsection{One independent scalar variable}

Our normal form is

\[
\frac{d}{dx}\left(f(x)\right).
\]

Read it as “differentiate the parenthesized expression with respect to \(x\).” The derivative operator is not enclosed in another pair of parentheses. No prime marks, time dots, bare function names as operands, or implicit differentiation variables appear in worked calculations.

For second derivatives, use the established higher-order operator and expand it where introduced:

\[
\frac{d^2}{dx^2}\left(f(x)\right)
=
\frac{d}{dx}\left(\frac{d}{dx}\left(f(x)\right)\right).
\]

The superscript \(2\) means “differentiate twice,” not “square the derivative.” Every worked higher-derivative calculation shows the intermediate first derivative. This is clearer than requiring ever-deeper nesting in every formula.

Differentiate first and evaluate second:

\[
\left.\frac{d}{dx}\left(f(x)\right)\right|_{x=a}.
\]

Always calculate the derivative before substituting the specified input.

\subsection{Several independent scalar variables}

Use the established partial-derivative operator and put the held-fixed instruction beside it:

\[
\frac{\partial}{\partial x}\left(f(x,y)\right),
\qquad y\text{ held fixed}.
\]

A nearby sentence lists all other independent inputs and fixed parameters. This avoids putting a paragraph into an operator subscript. If which quantity stays fixed changes the answer, state it on each affected line. A bar with a subscript is reserved for evaluation, so it is not also used to mean “hold this variable fixed.”

Keep ordinary \(d\) and partial \(\partial\) distinct. Partial differentiation changes one independent coordinate while holding the others fixed. Ordinary differentiation applies when the expression has one independent scalar input, including a whole path-dependent composition.

For a differentiable scalar-valued function \(f(x,y)\) and differentiable scalar paths \(x(t),y(t)\), write

\[
\begin{aligned}
\frac{d}{dt}\left(f(x(t),y(t))\right)
={}&\left.\frac{\partial}{\partial x}\left(f(x,y)\right)
\right|_{x=x(t),\,y=y(t)}
\frac{d}{dt}\left(x(t)\right)\\
&+\left.\frac{\partial}{\partial y}\left(f(x,y)\right)
\right|_{x=x(t),\,y=y(t)}
\frac{d}{dt}\left(y(t)\right).
\end{aligned}
\]

In the first partial derivative hold \(y\) fixed; in the second hold \(x\) fixed. Calculate each partial derivative with independent coordinates, evaluate it on the path, and only then multiply by the corresponding input rate.

Mixed partial derivatives are initially written as nested operations, with the order stated in words. Do not assume that reversing the order gives the same result without appropriate hypotheses.

\subsection{Vector and matrix outputs}

For a single scalar input, differentiate component by component in a fixed basis. For example, \(\mathbf r(t)\) is a two-component column vector, and \(r_1(t),r_2(t)\) are scalar functions:

\[
\mathbf r(t)=\begin{bmatrix}r_1(t)\\r_2(t)\end{bmatrix},
\qquad
\frac{d}{dt}\left(\mathbf r(t)\right)
=\begin{bmatrix}
\frac{d}{dt}\left(r_1(t)\right)\\
\frac{d}{dt}\left(r_2(t)\right)
\end{bmatrix}.
\]

This separates the round derivative-operand delimiters from the square vector-display delimiters. Matrix and tensor examples likewise declare the derivative's shape.

For vector inputs, introduce named scalar coordinates first. Let \(\mathbf F(x_1,\ldots,x_n)\in\mathbb R^m\) be differentiable. Its Jacobian is the \(m\times n\) matrix with scalar entries

\[
J_{ij}(x_1,\ldots,x_n)
=\frac{\partial}{\partial x_j}\left(F_i(x_1,\ldots,x_n)\right).
\]

All \(x_k\) with \(k\ne j\) are held fixed. Index \(i\) selects an output, \(1\leq i\leq m\); index \(j\) selects an input, \(1\leq j\leq n\). In concrete examples list the actual coordinates instead of using ellipses.

For a differentiable scalar function \(f(x,y)\) in ordinary Euclidean coordinates, define the gradient explicitly:

\[
\nabla f(x,y)=
\begin{bmatrix}
\frac{\partial}{\partial x}\left(f(x,y)\right)\\
\frac{\partial}{\partial y}\left(f(x,y)\right)
\end{bmatrix}.
\]

Hold the other coordinate fixed in each entry. The gradient is a column vector; the scalar function's Jacobian is its transpose, a one-row matrix. We do not use an unexplained fraction “with respect to a vector” for both objects. The superscript \(\mathsf T\) means transpose, never a derivative.

\subsection{Matrix inputs and tensor inputs}

Start with derivatives with respect to individual scalar entries. For a scalar function \(\phi(\mathbf A)\) of a real \(m\times n\) matrix with independently variable entries, the matrix gradient \(\nabla_{\mathbf A}\phi(\mathbf A)\) has the same shape as \(\mathbf A\). Its \((i,j)\) entry is the scalar partial derivative with respect to \(A_{ij}\), holding all other entries fixed. If the matrix is symmetric or otherwise constrained, identify its independent parameters before differentiating; its entries cannot all be treated as independent.

For a general differentiable matrix-valued function \(\mathbf F(\mathbf A)\), use the established derivative-map form \(\mathrm D\mathbf F(\mathbf A)(\mathbf H)\) after defining it: it is the derivative at input \(\mathbf A\), applied to the increment matrix \(\mathbf H\). Declare both matrices' shapes, the output shape, and all parameters. This is a linear map applied to an increment, not matrix division. The final parentheses replace the square argument brackets used in some references.

Tensor calculations specify the tensor space, basis, component indices, and index ranges. For geometric tensors, upper and lower indices have distinct roles and must be explained; no implicit summation is allowed. Fixed-basis component differentiation is introduced first. Moving bases require additional terms.

\section{4. Integrals}

Use the usual definite-integral form:

\[
\int_a^b f(x)\,dx.
\]

Immediately identify \(x\) as the scalar integration variable, with lower limit \(x=a\) and upper limit \(x=b\). The differential \(dx\) remains on every integral. Parentheses group a sum when needed; a simple integrand does not need an extra wrapper. We do not routinely write \(x=\) inside the bounds: the nearby variable-and-bounds statement supplies that clarity while keeping the expression familiar.

The symbol \(x\) inside a definite integral is a bound integration variable. If the bounds are fixed and there are no other inputs, the result is a scalar number, not a function of \(x\).

For an indefinite integral on an appropriate interval,

\[
\int f(x)\,dx=F(x)+C,
\qquad \frac{d}{dx}\left(F(x)\right)=f(x).
\]

Here \(C\) is an arbitrary scalar constant independent of \(x\). For vector and matrix outputs, the arbitrary constant has the output's shape.

With other inputs held fixed,

\[
\int f(x,y)\,dx=F(x,y)+C(y),\qquad y\text{ held fixed}.
\]

The integration constant may depend on the input that was not integrated over. Definite integration over \(x\) can likewise leave a function of \(y\).

Write endpoint evaluation without another set of square brackets:

\[
\left.F(x)\right|_{x=a}^{x=b}=F(b)-F(a).
\]

Always show the subtraction. For multiple integrals, specify the order in words and give each integral its own bounds and differential. For curve, surface, and volume integrals, specify domain, measure, parameterization when used, and orientation when relevant. Arc length \(ds\), surface area \(dS\), and a vector displacement are different objects, each introduced with a definition and component formula.

\section{5. Substitution without hidden steps}

Introduce scalar functions \(g(x)\) and \(h(u)\) and a new scalar coordinate \(u=g(x)\). Before the substitution, \(u\) depends on \(x\); in the transformed integral, \(u\) is the integration variable.

For continuously differentiable \(g\) on \([a,b]\) and continuous \(h\) on an interval containing its image,

\[
\int_a^b h(g(x))\frac{d}{dx}\left(g(x)\right)\,dx
=\int_{g(a)}^{g(b)}h(u)\,du.
\]

On the left, integrate over \(x\) from \(a\) to \(b\). On the right, integrate over \(u\) from \(g(a)\) to \(g(b)\). This form does not require an inverse for \(g\). A substitution using an inverse may need domain restrictions and a branch choice.

Required steps: define the substitution; compute its derivative; match every factor; transform both bounds; write the complete new integral; evaluate; check. For indefinite integration, substitute back and verify by differentiation.

Differential notation such as \(du=2x\,dx\) is conventional, not inherently invalid. We teach its meaning before using it: when \(u=g(x)=x^2+1\), the differential of \(u\) is the derivative of \(g(x)\) multiplied by the differential of \(x\). In early worked solutions the full substitution formula remains visible. Never present canceling the letters \(d\) and \(x\) as a proof, and never divide by a possibly zero derivative without justification.

\section{6. Other notation collisions to prevent}

\begin{itemize}
\item \textbf{Finite change versus differential:} \(\Delta x\) is a finite scalar increment. A differential gives the linear part of a change; it is not automatically the exact finite change. Introduce \(\mathrm Df(x,y)(\Delta x,\Delta y)\) by its full partial-derivative formula before using it.
\item \textbf{Function inverse versus reciprocal:} write \(1/f(x)\) for a reciprocal and \(f^{-1}(y)\) only for a defined inverse function. Use \(\arcsin(x)\) rather than \(\sin^{-1}(x)\).
\item \textbf{Function powers:} write \((\sin(x))^2\), not the compressed \(\sin^2 x\). Every elementary function shows its input.
\item \textbf{Matrix inverse versus reciprocal:} \(\mathbf A^{-1}\) is a matrix inverse when it exists; entrywise reciprocals are defined separately. Do not use matrix “division.”
\item \textbf{Absolute value versus determinant:} use \(|s|\) for a scalar absolute value, \(\det(\mathbf A)\) for a determinant, and \(\lVert\mathbf v\rVert\) for a vector norm. Define the chosen norm.
\item \textbf{Products:} scalar multiplication, vector dot products, cross products, matrix products, tensor products, and contractions are named when introduced. Matrix factors retain their order. Entrywise products are explicitly identified.
\item \textbf{Summations:} write every summation sign and its limits. State free-index ranges. Expand all terms in first small examples. Einstein summation is explained in the translation guide only.
\item \textbf{Tensor meaning:} multilinear tensors and multidimensional data arrays are related through chosen representations but are not silently treated as identical concepts.
\item \textbf{Equality versus approximation:} use \(=\) for exact equality and \(\approx\) for approximation. State the approximation's conditions and error estimate, or explicitly say that no error bound has been established.

\end{itemize}
\section{7. Worked-solution presentation}

Every solution uses numbered steps and textual labels:

\begin{longtable}{@{}>{\raggedright\arraybackslash}p{0.300\textwidth}>{\raggedright\arraybackslash}p{0.300\textwidth}>{\raggedright\arraybackslash}p{0.300\textwidth}@{}}
Label & Meaning & Planned print styling \\[5pt]
\hline\endhead
SETUP & Types, dependencies, domain, target & Normal text \\[5pt]
CHOICE & Why this method applies; cues to recognize it & Normal text \\[5pt]
CALCULUS & Derivative or integral rule being applied & Black math with a named rule \\[5pt]
ALGEBRA & Expansion, factoring, cancellation, solving, rearrangement & Dark blue \texttt{\#244A73} with an ALGEBRA label \\[5pt]
IDENTITY & Trigonometric, exponential, logarithmic, or other exact identity & Same dark blue with an IDENTITY label \\[5pt]
TRICK & A supporting identity, rewrite, or approximation tool from the tricks appendix & Purple \texttt{\#6F398E}, with a TRICK label and an appendix reference \\[5pt]
APPROXIMATION & A deliberately approximate replacement & Same dark blue with an APPROXIMATION label and conditions \\[5pt]
CHECK & Verification, dimensions, units, domain, or plausibility & Normal text \\[5pt]
\end{longtable}

Color is supplementary: the labels survive grayscale printing and text-only tutoring. Do not make algebra light gray or low contrast. Markdown pilot files use the labels; the typeset book will apply the color styles centrally.

Separate calculus from algebra into distinct lines. Give the unsimplified result before simplifying. Explain each factor, sign, and canceled term. State conditions for division and cancellation. Name an identity and show its general form before substituting the problem's expressions.

When a step invokes a supporting tool from the tricks appendix,
TRICK styling takes precedence over IDENTITY or APPROXIMATION styling. Supporting
tools in that appendix are purple; algebra in the main chapters remains dark blue.
Substitution and integration by parts are main calculus methods and are not labeled TRICK.

Each problem has a stable ID, a problem-only prompt, a graduated hint ladder, a separate complete solution, a verification, a method-recognition note, and a nearby practice variant. Solutions must not sit beside the question in the final page layout.

\section{8. Tutor contract}

When asked to help with a problem, the tutor uses its ID and this standard. Begin at the learner's last understood step, ask for one manageable next action, and wait. Explain a needed prerequisite locally. Do not skip algebra or switch notation to save space. Give the full solution when requested; otherwise reveal only the next requested hint or step. A valid alternative method is welcome if its assumptions and dependencies are equally explicit.

\section{9. Editorial acceptance checklist}

\begin{itemize}
\item Can the reader identify every object's type and every function input locally?
\item Is the differentiation/integration variable explicit, including what is held fixed?
\item Are intermediate dependencies, bounds, domains, and arbitrary constants correct?
\item Is each equation exact, or explicitly marked approximate?
\item Are calculus rules distinguished from algebra and identities?
\item Can every transition be reproduced without a hidden manipulation?
\item Is there a check appropriate to the problem and a cue explaining the method choice?
\item Do the question, hints, and solution share one stable problem ID?
\end{itemize}

\chapter{The curriculum}

The book is a review and problem-solving manual. Short concept refreshers support worked problems; readers are not expected to remember a formula or an algebraic manipulation merely because they once took a class.

The two main parts are \textbf{Differentiation and differentials} and \textbf{Integration}. Within each part, complexity grows from scalar variables and outputs to vector, matrix, and tensor objects. Input dimension and integration geometry form a second organizing axis: an object's type alone does not determine the calculus method.

\section{Entry toolkit: how to read and use the book}

\begin{enumerate}
\item The notation standard, dependency diagrams, types, shapes, domains, and units.
\item A diagnostic set for fractions, powers, roots, logarithms, exponentials, factoring, completing the square, equations, inequalities, and trigonometry.
\item A referenced toolbox of exact identities and approximation rules. Each entry has conditions, a worked example, and links to problems using it.
\item Essential reminders about limits, continuity, and what a derivative or integral asks for. Revisit limits later where a method needs them.

\end{enumerate}
Diagnostic results direct readers to local refreshers; they do not block access to calculus chapters.

\section{Part I: Differentiation and differentials}

\begin{longtable}{@{}>{\raggedright\arraybackslash}p{0.300\textwidth}>{\raggedright\arraybackslash}p{0.300\textwidth}>{\raggedright\arraybackslash}p{0.300\textwidth}@{}}
Chapter & Scope & Problem-solving focus \\[5pt]
\hline\endhead
D1 & Scalar input, scalar output & Read dependencies; use elementary, sum, and constant-multiple rules; check from the definition in selected examples \\[5pt]
D2 & Products, quotients, and compositions & Recognize structure; choose whether to simplify first; fully expose the chain rule \\[5pt]
D3 & Implicit, inverse, and parametric relationships & Decide what depends on what; solve for a requested rate; handle logarithmic differentiation and domain restrictions \\[5pt]
D4 & Higher derivatives, local approximation, and optimization & Repeated operations, linearization, Taylor polynomials with errors, critical points, and endpoint checks \\[5pt]
D5 & Several scalar inputs, scalar output & Partial derivatives, gradients, directional derivatives, Hessians, tangent approximations, and constrained optimization \\[5pt]
D6 & Vector outputs and vector inputs & Velocity and acceleration; Jacobians; multivariable compositions; derivative as a linear map \\[5pt]
D7 & Matrix inputs and outputs & Entrywise time derivatives; products and inverses; scalar functions of matrices; explicit index conventions and shape checks \\[5pt]
D8 & Tensor-valued functions and contractions & Fixed-basis component derivatives; multilinear structure; tensor products and contractions; careful separation of tensor meaning from array storage \\[5pt]
\end{longtable}

An optional extension follows D8: changing bases, curvilinear coordinates, and covariant derivatives. This is a separate prerequisite-bearing topic, not a silent extension of componentwise differentiation.

\section{Part II: Integration}

\begin{longtable}{@{}>{\raggedright\arraybackslash}p{0.300\textwidth}>{\raggedright\arraybackslash}p{0.300\textwidth}>{\raggedright\arraybackslash}p{0.300\textwidth}@{}}
Chapter & Scope & Problem-solving focus \\[5pt]
\hline\endhead
I1 & Scalar antiderivatives and definite integrals & Meaning, fundamental theorem, bounds, constants, signed accumulation, and basic rules \\[5pt]
I2 & Change of variable in one dimension & Recognize a composition and its derivative; transform bounds; handle inverse branches and domain restrictions \\[5pt]
I3 & Integration by parts & Choose factors; expose product-rule origins; repeat operations and solve for a recurring integral \\[5pt]
I4 & Algebraic and trigonometric methods & Polynomial division, partial fractions, completing the square, trig identities, and trig substitution \\[5pt]
I5 & Choosing a method for unfamiliar integrals & Mixed practice without method labels; compare valid routes; recognize when elementary antiderivatives are unavailable \\[5pt]
I6 & Improper integrals and numerical approximation & Limits at singularities and infinity; convergence before evaluation; quadrature and error bounds \\[5pt]
I7 & Vector, matrix, and tensor outputs over a scalar variable & Component integrals; initial conditions; shape preservation; fixed-basis assumptions and limits of componentwise reasoning \\[5pt]
I8 & Double and triple integrals & Regions, bounds, order changes, signed and nonnegative integrands, and conditions for interchanging integrals \\[5pt]
I9 & Multivariable changes of coordinates & Jacobian determinants, absolute values, transformed regions, polar/cylindrical/spherical examples \\[5pt]
I10 & Curves: scalar accumulation and vector work & Parameterizations, arc-length measure, tangents, orientation, and distinguishing two kinds of line integral \\[5pt]
I11 & Surfaces: scalar accumulation and vector flux & Surface measure, normal direction, parameterization, and orientation \\[5pt]
I12 & Connecting local derivatives with global integrals & Green's, divergence, and Stokes' theorems; choose the easier side and check hypotheses \\[5pt]
I13 & Integrals depending on parameters and mixed objects & Differentiate under an integral with stated conditions; vector/matrix/tensor integrands over regions; explicit contractions \\[5pt]
\end{longtable}

Optional advanced extensions: differential forms and generalized Stokes; integration on manifolds; matrix exponentials and systems of differential equations; higher-order tensors in applications. These are scoped separately so the core review remains usable.

\section{Dependencies and reading routes}

\begin{itemize}
\item The quickest scalar review is D1--D4 followed by I1--I6.
\item I2 uses the chain rule from D2; I3 uses the product rule from D2.
\item I7 can follow the scalar integration chapters once component notation is understood.
\item I8 needs scalar integration and elementary multivariable functions; I9 also needs the Jacobian from D6 and determinant algebra.
\item I10--I12 draw on D5--D6 and the relevant integration chapters. Introduce divergence and curl explicitly before their theorems.
\item D7--D8 and the tensor portions of I7/I13 have short linear-algebra prerequisite modules. Their unfamiliar concepts are taught locally rather than assumed.

\end{itemize}
\section{Repeated chapter structure}

\begin{enumerate}
\item A short concept refresher and a “what can vary?” map.
\item A rule sheet using only the book's notation, with conditions and output types.
\item Fully worked anchor problems, starting with a single new idea.
\item Near-transfer practice: a changed sign, exponent, bound, dependency, or dimension.
\item Mixed problems with the method name withheld.
\item Failure-analysis problems: diagnose a plausible wrong step.
\item A method-selection page: cues, necessary checks, and common dead ends.
\item Hints and full solutions in separate, clearly indexed sections.

\end{enumerate}
Every practice problem will have a complete solution. Every approximation will identify its regime of validity. Checks include reverse differentiation for antiderivatives, shape and unit checks where applicable, and independent evaluations when useful.

The eventual book should provide both a printable reading experience and stable problem/step references suitable for an LLM tutor. First settle the notation using the pilots, then develop one representative scalar chapter before expanding to the remaining chapters.

\chapter{D1: Scalar differentiation foundations}

Alpha chapter. This chapter develops powers, sums, fixed parameters, domains, and
evaluation. Exponential, logarithmic, and trigonometric derivatives will extend
this foundation in later material. Every problem below has hints, a complete
solution, and a check.

\section{What this chapter trains}

Before reaching for a rule, identify the independent variable and inspect the
expression. Sometimes rewriting the algebra makes the calculus much simpler.
Your goal is to produce both a derivative and a statement of where it is valid.

All functions here have scalar inputs and scalar outputs. A derivative with
respect to a scalar input is a scalar-valued function; evaluating it at a
specified input gives a scalar number.

\subsection{A repeatable approach}

\begin{enumerate}
\item \textbf{Declare:} identify the function, its inputs, fixed parameters, and domain.
\item \textbf{Inspect:} separate sums; distinguish fixed coefficients from changing expressions.
\item \textbf{Rewrite if useful:} turn roots and reciprocals into powers, or simplify a fraction without losing its original domain.
\item \textbf{Differentiate:} name each rule and retain the function arguments.
\item \textbf{Simplify:} show the coefficient and exponent arithmetic.
\item \textbf{Evaluate if requested:} substitute only after calculating the derivative.
\item \textbf{Check:} compare with the definition, check a sign, or use another valid calculation. State what that check actually establishes.

\end{enumerate}
\subsection{When these rules are sufficient}

A sum of constant multiples of powers of the independent variable fits this
chapter. A changing expression raised to a power, such as \((x^2+1)^4\), requires
the chain rule introduced in D2-001. A product of changing
factors or a quotient that does not simplify may require the product or quotient
rule. Never differentiate the factors of a product separately and multiply the
results.

\section{Rules with their assumptions}

In this rule sheet \(x\) is an independent real scalar, \(c\) and \(p\) are fixed real
scalar parameters, and \(u(x),v(x)\) are scalar functions differentiable at the
inputs under consideration. The output of each derivative is scalar.

\subsection{Constant and constant-multiple rules}

\[
\frac{d}{dx}\left(c\right)=0,
\qquad
\frac{d}{dx}\left(cu(x)\right)
=c\frac{d}{dx}\left(u(x)\right).
\]

“Constant” means independent of the chosen variable. A letter such as \(a\) can be
a constant with respect to \(x\) even though different problems assign it different
values. A function such as \(a(x)\) is not a fixed coefficient.

\subsection{Sum and difference rules}

\[
\frac{d}{dx}\left(u(x)+v(x)\right)
=\frac{d}{dx}\left(u(x)\right)
+\frac{d}{dx}\left(v(x)\right).
\]

For a difference, retain the minus sign between the derivatives. Apply these rules
repeatedly to longer sums, writing each term explicitly.

\subsection{Power rule}

For a fixed real exponent \(p\) and \(x>0\),

\[
\frac{d}{dx}\left(x^p\right)=p x^{p-1}.
\]

Nonnegative integer powers are differentiable for all real \(x\); handle the
constant power \(x^0=1\) directly with the constant rule. Negative integer powers
require \(x\ne0\). Fractional powers need their own real-domain check. The positive
domain version avoids ambiguous real powers of negative inputs, but it does not
say that every rational power is restricted to positive inputs.

This rule concerns \(x^p\) with a fixed exponent. It does not directly apply to
\(p^x\), where the exponent changes, or to a power of an inner function without
accounting for that function's derivative.

\subsection{Definition for checking a derivative}

Declare \(h\) as a nonzero real scalar increment. At an interior input \(x\) where
the two-sided limit exists,

\[
\frac{d}{dx}\left(f(x)\right)
=\lim_{h\to0}\frac{f(x+h)-f(x)}{h}.
\]

Require \(x+h\) to lie in the function's domain. You may divide by \(h\) before
taking the limit because the quotient uses nonzero increments. Substituting
\(h=0\) into that quotient before simplifying would divide by zero.

These are standard derivative definitions and rules; see
\href{https://openstax.org/books/calculus-volume-1/pages/3-1-defining-the-derivative}{OpenStax, Calculus Volume 1, §3.1}
and \href{https://openstax.org/books/calculus-volume-1/pages/3-3-differentiation-rules}{§3.3}.
The problems and worked explanations below are original to this workbook.

\section{Problems}

\subsection{D1-001 --- A polynomial, term by term}

\textbf{Objects and dependencies.} \(x\in\mathbb R\) is the independent scalar variable.
The scalar function is \(f(x)=4x^3-5x^2+7x-9\), defined for every real \(x\).
All coefficients are fixed real scalars. Find the scalar-valued function

\[
\frac{d}{dx}\left(f(x)\right).
\]

\subsection{D1-002 --- A root and a reciprocal}

\textbf{Objects and dependencies.} \(x\) is an independent real scalar with \(x>0\).
The scalar function is

\[
f(x)=5\sqrt{x}-\frac{3}{x^2}+2.
\]

The square root is the nonnegative real square root. All coefficients are fixed
scalars. Find the derivative with respect to \(x\) and evaluate it at \(x=4\).

\subsection{D1-003 --- Letters that stay fixed}

\textbf{Objects and dependencies.} \(x\in\mathbb R\) is the independent scalar variable.
\(a,b\in\mathbb R\) are fixed scalar parameters. The scalar output is

\[
f(x;a,b)=ax^2+bx+4.
\]

Differentiate with respect to \(x\), holding \(a\) and \(b\) fixed. Then evaluate the
derivative at \(x=3\) with parameter values \(a=2\) and \(b=-1\).

\subsection{D1-004 --- Simplification with a missing input}

\textbf{Objects and dependencies.} \(x\) is an independent real scalar with \(x\ne3\).
The scalar function is

\[
f(x)=\frac{x^2-9}{x-3}.
\]

Find its derivative on its domain. A student cancels a factor and concludes that
this original function has derivative \(1\) at \(x=3\) too. Explain whether that
conclusion follows.

\subsection{D1-005 --- Build a derivative from its definition}

\textbf{Objects and dependencies.} \(x\in\mathbb R\) is the independent scalar variable.
The scalar function is \(f(x)=x^2+2x\). Use the limit definition to calculate its
derivative, without using the power rule as the derivation. Then evaluate the
derivative at \(x=-1\).

During the calculation, \(h\in\mathbb R\) is a nonzero scalar increment that tends
to zero; it is not a new input of \(f\).

\subsection{D1-006 --- Choose a method for a rational expression}

\textbf{Objects and dependencies.} \(x\) is an independent real scalar with \(x\ne0\).
The scalar function is

\[
f(x)=\frac{2x^3-3x+4}{x}.
\]

Find the scalar-valued derivative and the scalar slope at \(x=2\).
Choose and explain your method before calculating.

\section{Hint ladders}

\subsection{D1-001}

\begin{enumerate}
\item Identify the four terms and keep the signs attached to them.
\item The fixed coefficients can be taken outside their derivative operators; the constant term has derivative zero.
\item Differentiate \(x^3\), \(x^2\), and \(x\) separately, then multiply their coefficients and simplify their exponents.

\end{enumerate}
\subsection{D1-002}

\begin{enumerate}
\item Rewrite both the root and the reciprocal as powers of the same independent variable.
\item On \(x>0\), use \(\sqrt{x}=x^{1/2}\) and \(1/x^2=x^{-2}\).
\item The derivative of the negative reciprocal term has a positive coefficient because \((-3)(-2)=6\). Evaluate at \(x=4\) after differentiating.

\end{enumerate}
\subsection{D1-003}

\begin{enumerate}
\item The symbols \(a\) and \(b\) are fixed with respect to \(x\); they are not functions \(a(x)\) and \(b(x)\).
\item Treat \(a\) just as you would treat the coefficient \(4\) in D1-001; treat \(b\) just as you would treat \(7\).
\item Obtain the derivative with its parameters visible, then substitute \(x=3\), \(a=2\), and \(b=-1\).

\end{enumerate}
\subsection{D1-004}

\begin{enumerate}
\item Can the numerator be factored using a difference-of-squares identity?
\item Write \(x^2-9=(x-3)(x+3)\). Cancellation is allowed only when \(x-3\ne0\).
\item The simplified formula describes the original function only on \(x\ne3\). A formula that fills in the missing value defines a different function.

\end{enumerate}
\subsection{D1-005}

\begin{enumerate}
\item Write \(f(x+h)\) using the same function definition, replacing every input \(x\) by \(x+h\).
\item Expand the square, then subtract all of \(f(x)\) inside parentheses.
\item Factor a single \(h\) from the numerator, divide for \(h\ne0\), and only then take the limit.

\end{enumerate}
\subsection{D1-006}

\begin{enumerate}
\item Inspect the denominator: it is a single power of the independent variable.
\item Divide each numerator term by \(x\) separately while preserving the condition \(x\ne0\).
\item The resulting expression is a sum of powers. Calculate the derivative, then evaluate its value at \(x=2\).

\end{enumerate}
\clearpage
\section{Complete solutions}

\subsection{D1-001}

\textbf{Step 1 --- SETUP.} \(x\in\mathbb R\) is independent and scalar; \(f(x)=4x^3-5x^2+7x-9\)
is scalar. The coefficients do not vary with \(x\). A polynomial is differentiable
at every real input.

\textbf{Step 2 --- CHOICE.} This is a sum of constant multiples of powers of \(x\).
Use sum, difference, constant-multiple, power, and constant rules.

\textbf{Step 3 --- CALCULUS: separate the terms.}

\[
\begin{aligned}
\frac{d}{dx}\left(f(x)\right)
&=\frac{d}{dx}\left(4x^3-5x^2+7x-9\right)\\
&=4\frac{d}{dx}\left(x^3\right)
-5\frac{d}{dx}\left(x^2\right)
+7\frac{d}{dx}\left(x\right)
-\frac{d}{dx}\left(9\right)\\
&=4(3x^{3-1})-5(2x^{2-1})+7(1)-0.
\end{aligned}
\]

\begingroup\color{AlgebraBlue}
\textbf{Step 4 --- ALGEBRA: simplify exponents and coefficients.}

\[
\begin{aligned}
3-1&=2,\qquad 2-1=1,\\
4(3x^2)-5(2x)+7(1)-0
&=(4\cdot3)x^2-(5\cdot2)x+7\\
&=12x^2-10x+7.
\end{aligned}
\]

\endgroup
\textbf{Result.}

\[
\boxed{\frac{d}{dx}\left(f(x)\right)=12x^2-10x+7,
\qquad x\in\mathbb R.}
\]

\textbf{Step 5 --- CHECK: compare at one input with the definition.} At \(x=0\), the
formula gives \(12(0)^2-10(0)+7=7\). For a nonzero scalar increment \(h\),
\(f(h)=4h^3-5h^2+7h-9\) and \(f(0)=-9\).

\begingroup\color{AlgebraBlue}
\textbf{Step 6 --- ALGEBRA: form the check's difference quotient.}

\[
\begin{aligned}
\frac{f(h)-f(0)}{h}
&=\frac{4h^3-5h^2+7h-9-(-9)}{h}\\
&=\frac{4h^3-5h^2+7h}{h}\\
&=\frac{h(4h^2-5h+7)}{h}\\
&=4h^2-5h+7,\qquad h\ne0.
\end{aligned}
\]

\endgroup
\textbf{Step 7 --- CHECK: take the limit.} As \(h\to0\), the expression tends to \(7\),
matching the derivative at zero. This checks one input; it is not an independent
proof of the formula at every input.

\textbf{Recognition cue.} Separate a sum term by term, but keep the signs. A fixed
coefficient stays as a multiplier; a standalone constant contributes zero.

\subsection{D1-002}

\textbf{Step 1 --- SETUP.} \(x>0\) is the independent real scalar.
\(f(x)=5\sqrt{x}-3/x^2+2\) is scalar, with fixed coefficients. The positive domain
both defines the root smoothly and excludes division by zero.

\textbf{Step 2 --- CHOICE.} Rewriting the root and reciprocal as powers makes the power
rule available for both changing terms.

\begingroup\color{AlgebraBlue}
\textbf{Step 3 --- ALGEBRA: rewrite the function on its domain.}

\[
\sqrt{x}=x^{1/2},\qquad \frac{1}{x^2}=x^{-2},\qquad
f(x)=5x^{1/2}-3x^{-2}+2,\qquad x>0.
\]

\endgroup
\textbf{Step 4 --- CALCULUS: apply the rules before simplifying.}

\[
\begin{aligned}
\frac{d}{dx}\left(f(x)\right)
&=5\frac{d}{dx}\left(x^{1/2}\right)
-3\frac{d}{dx}\left(x^{-2}\right)
+\frac{d}{dx}\left(2\right)\\
&=5\left(\frac12 x^{1/2-1}\right)
-3\left((-2)x^{-2-1}\right)+0.
\end{aligned}
\]

\begingroup\color{AlgebraBlue}
\textbf{Step 5 --- ALGEBRA: handle the signs and powers.}

\[
\begin{aligned}
\frac12-1&=\frac12-\frac22=-\frac12,\\
-2-1&=-3,\qquad (-3)(-2)=6,\\
5\left(\frac12 x^{-1/2}\right)-3\left((-2)x^{-3}\right)+0
&=\frac52 x^{-1/2}+6x^{-3}\\
&=\frac{5}{2\sqrt{x}}+\frac{6}{x^3}.
\end{aligned}
\]

\endgroup
\textbf{Step 6 --- CALCULUS: evaluate the already calculated derivative.}

\[
\left.\frac{d}{dx}\left(f(x)\right)\right|_{x=4}
=\frac{5}{2\sqrt{4}}+\frac{6}{4^3}.
\]

\begingroup\color{AlgebraBlue}
\textbf{Step 7 --- ALGEBRA: combine the fractions.}

\[
\begin{aligned}
\sqrt4&=2,\qquad 4^3=4\cdot4\cdot4=64,\\
\frac{5}{2(2)}+\frac6{64}
&=\frac54+\frac3{32}\\
&=\frac{5\cdot8}{4\cdot8}+\frac3{32}\\
&=\frac{40+3}{32}=\frac{43}{32}.
\end{aligned}
\]

\endgroup
\textbf{Result.} The derivative is \(5/(2\sqrt{x})+6/x^3\) for \(x>0\), and its value at
\(x=4\) is the scalar number \(43/32\).

\textbf{Step 8 --- CHECK: sign and domain.} Both derivative terms are positive for
\(x>0\). This agrees with the behavior of the original terms: \(5\sqrt{x}\) increases,
and \(-3/x^2\) becomes less negative as \(x\) increases. The derivative formula also
retains the restriction \(x>0\). This sign check catches a likely sign mistake; it
does not establish the exact magnitude by itself.

\textbf{Recognition cue.} Convert roots and reciprocal powers before calculating.
Subtracting a term whose power-rule coefficient is negative produces a positive
term; write the multiplication to see why.

\subsection{D1-003}

\textbf{Step 1 --- SETUP.} \(x\in\mathbb R\) is the independent scalar.
\(a,b\in\mathbb R\) are fixed scalar parameters. The output \(f(x;a,b)=ax^2+bx+4\)
is scalar. Its derivative with respect to \(x\) is taken with both parameters fixed.

\textbf{Step 2 --- CHOICE.} Apply the same constant-multiple rules as for numerical
coefficients. A parameter's value is not required to perform this derivative.

\textbf{Step 3 --- CALCULUS.} Holding \(a\) and \(b\) fixed,

\[
\begin{aligned}
\frac{d}{dx}\left(f(x;a,b)\right)
&=a\frac{d}{dx}\left(x^2\right)
+b\frac{d}{dx}\left(x\right)
+\frac{d}{dx}\left(4\right)\\
&=a(2x)+b(1)+0.
\end{aligned}
\]

\begingroup\color{AlgebraBlue}
\textbf{Step 4 --- ALGEBRA.}

\[
a(2x)+b(1)+0=2ax+b.
\]

\endgroup
\textbf{Step 5 --- CALCULUS: substitute the input and parameter values.}

\[
\left.\frac{d}{dx}\left(f(x;a,b)\right)
\right|_{x=3,\,a=2,\,b=-1}
=2(2)(3)+(-1).
\]

\begingroup\color{AlgebraBlue}
\textbf{Step 6 --- ALGEBRA.}

\[
2(2)(3)+(-1)=4\cdot3-1=12-1=11.
\]

\endgroup
\textbf{Result.} The derivative is \(2ax+b\), with \(a,b\) fixed, for all real \(x\).
The requested evaluated slope is the scalar number \(11\).

\textbf{Step 7 --- CHECK: differentiate after fixing the same parameters.} First define
the scalar function \(g(x)=f(x;2,-1)=2x^2-x+4\). Its derivative is

\[
\frac{d}{dx}\left(g(x)\right)=2(2x)-1+0=4x-1.
\]

The coefficient simplification is \textbf{ALGEBRA}. At \(x=3\), \(4(3)-1=12-1=11\),
matching the parameterized calculation. This consistency check applies because
the parameters really are fixed. If \(a\) were replaced by a function \(a(x)\), the
dependencies and required rules would change.

\textbf{Recognition cue.} “Constant” is a relationship to the variable of
differentiation, not a requirement that the symbol be a numeral.

\subsection{D1-004}

\textbf{Step 1 --- SETUP.} \(x\ne3\) is the independent real scalar and
\(f(x)=(x^2-9)/(x-3)\) is scalar. The original definition excludes \(x=3\).

\textbf{Step 2 --- CHOICE.} Factor before choosing a quotient rule. The numerator has
a difference-of-squares form.

\begingroup\color{TrickPurple}
\textbf{Step 3 --- TRICK: difference of squares.} Use the identity in
T-005. For scalar expressions \(r,s\),

\[
r^2-s^2=(r-s)(r+s).
\]

With \(r=x\) and \(s=3\),

\[
x^2-9=x^2-3^2=(x-3)(x+3).
\]

\endgroup
\begingroup\color{AlgebraBlue}
\textbf{Step 4 --- ALGEBRA: cancel only on the original domain.}

\[
\begin{aligned}
f(x)&=\frac{(x-3)(x+3)}{x-3}\\
&=\frac{x-3}{x-3}(x+3)\\
&=1(x+3)=x+3,\qquad x\ne3.
\end{aligned}
\]

\endgroup
\textbf{Step 5 --- CALCULUS: differentiate the equivalent local formula.} At each
allowed input \(x\ne3\), the function agrees with \(x+3\) in a neighborhood:

\[
\frac{d}{dx}\left(f(x)\right)
=\frac{d}{dx}\left(x+3\right)
=\frac{d}{dx}\left(x\right)+\frac{d}{dx}\left(3\right)
=1+0=1.
\]

The last addition is \textbf{ALGEBRA}.

\textbf{Result.} The derivative is \(1\) for \(x\ne3\). The original function has no
derivative at \(x=3\) because it has no value there. Defining a new scalar function
\(\widetilde f(x)=x+3\) for every real \(x\) fills that hole; this extension has a
derivative at \(3\), but it is a different function with a different domain.

\textbf{Step 6 --- CHECK: keep the domain in the difference quotient.} For an allowed
input \(x\ne3\), take a nonzero scalar increment \(h\) small enough that \(x+h\ne3\).

\begingroup\color{AlgebraBlue}
\textbf{Step 7 --- ALGEBRA.}

\[
\frac{f(x+h)-f(x)}{h}
=\frac{(x+h+3)-(x+3)}{h}
=\frac{x+h+3-x-3}{h}
=\frac{h}{h}=1.
\]

\endgroup
\textbf{Step 8 --- CHECK.} The quotient's limit is \(1\) at every allowed input.
This verifies the derivative on the original domain and supplies no derivative
at the excluded input. The student's claim confuses the function with its extension.

\textbf{Recognition cue.} Simplification may shorten a formula while leaving a hole
in its domain. Carry the domain alongside the simplified expression.

\subsection{D1-005}

\textbf{Step 1 --- SETUP.} \(x\in\mathbb R\) is independent, \(f(x)=x^2+2x\) is scalar,
and \(h\ne0\) is a real scalar increment tending to zero. All inputs \(x+h\) are
allowed because the function is defined on all of \(\mathbb R\).

\textbf{Step 2 --- CHOICE.} Use the limit definition. The order is substitution,
expansion, subtraction, division, and then the limit.

\textbf{Step 3 --- CALCULUS: write the definition with the full function.}

\[
\frac{d}{dx}\left(f(x)\right)
=\lim_{h\to0}\frac{f(x+h)-f(x)}{h}.
\]

\begingroup\color{TrickPurple}
\textbf{Step 4 --- TRICK: expand the shifted square.} Use
T-003. The scalar identity
\((r+s)^2=r^2+2rs+s^2\), with \(r=x\) and \(s=h\), gives

\[
f(x+h)=(x+h)^2+2(x+h)=x^2+2xh+h^2+2x+2h.
\]

The distribution \(2(x+h)=2x+2h\) is an \textbf{ALGEBRA} step.

\endgroup
\begingroup\color{AlgebraBlue}
\textbf{Step 5 --- ALGEBRA: subtract the entire original value.}

\[
\begin{aligned}
f(x+h)-f(x)
&=(x^2+2xh+h^2+2x+2h)-(x^2+2x)\\
&=x^2+2xh+h^2+2x+2h-x^2-2x\\
&=(x^2-x^2)+(2x-2x)+2xh+h^2+2h\\
&=2xh+h^2+2h\\
&=h(2x+h+2).
\end{aligned}
\]

\endgroup
\begingroup\color{AlgebraBlue}
\textbf{Step 6 --- ALGEBRA: divide before letting the increment tend to zero.}

\[
\frac{f(x+h)-f(x)}{h}
=\frac{h(2x+h+2)}{h}
=2x+h+2,\qquad h\ne0.
\]

\endgroup
\textbf{Step 7 --- CALCULUS: take the limit and then evaluate the derivative.} In the
limit, \(x\) stays fixed and \(h\to0\):

\[
\frac{d}{dx}\left(f(x)\right)=\lim_{h\to0}(2x+h+2)=2x+2.
\]

\[
\left.\frac{d}{dx}\left(f(x)\right)\right|_{x=-1}=2(-1)+2.
\]

\begingroup\color{AlgebraBlue}
\textbf{Step 8 --- ALGEBRA.}

\[
2(-1)+2=-2+2=0.
\]

\endgroup
\textbf{Result.} The derivative is \(2x+2\) for every real \(x\); the slope at \(x=-1\)
is the scalar number \(0\). The value \(f(-1)=1-2=-1\) is a function value, a
different quantity from its slope.

\textbf{Step 9 --- CHECK: compare with the power-rule calculation.}

\[
\frac{d}{dx}\left(x^2+2x\right)
=\frac{d}{dx}\left(x^2\right)+2\frac{d}{dx}\left(x\right)
=2x+2(1)=2x+2.
\]

The coefficient simplification is \textbf{ALGEBRA}. This agrees with the definition
calculation. Replacing \(x\) by \(-1\) before differentiating would give the constant
expression \(-1\), whose zero derivative would not tell you the general derivative
or justify the slope calculation. Here it would accidentally match the slope.

\textbf{Recognition cue.} Subtract the entire old function value in parentheses.
Cancellation of the increment is allowed only while it is nonzero.

\subsection{D1-006}

\textbf{Step 1 --- SETUP.} \(x\ne0\) is the independent real scalar;
\(f(x)=(2x^3-3x+4)/x\) is scalar. The requested slope at \(x=2\) is valid because
\(2\ne0\). There are no varying parameters.

\textbf{Step 2 --- CHOICE.} Every numerator term can be divided by the simple
denominator \(x\). Rewrite as powers and use sum and power rules. This avoids
introducing an unnecessary quotient-rule calculation.

\begingroup\color{AlgebraBlue}
\textbf{Step 3 --- ALGEBRA: divide each term without discarding the restriction.}

\[
\begin{aligned}
f(x)&=\frac{2x^3}{x}-\frac{3x}{x}+\frac4x\\
&=2x^{3-1}-3+4x^{-1}\\
&=2x^2-3+4x^{-1},\qquad x\ne0.
\end{aligned}
\]

\endgroup
\textbf{Step 4 --- CALCULUS: differentiate the rewritten function.}

\[
\begin{aligned}
\frac{d}{dx}\left(f(x)\right)
&=2\frac{d}{dx}\left(x^2\right)
-\frac{d}{dx}\left(3\right)
+4\frac{d}{dx}\left(x^{-1}\right)\\
&=2(2x)-0+4((-1)x^{-1-1}).
\end{aligned}
\]

\begingroup\color{AlgebraBlue}
\textbf{Step 5 --- ALGEBRA.}

\[
\begin{aligned}
-1-1&=-2,\\
2(2x)-0+4((-1)x^{-2})
&=4x-4x^{-2}=4x-\frac4{x^2}.
\end{aligned}
\]

\endgroup
\textbf{Step 6 --- CALCULUS: evaluate at the specified input.}

\[
\left.\frac{d}{dx}\left(f(x)\right)\right|_{x=2}
=4(2)-\frac4{2^2}.
\]

\begingroup\color{AlgebraBlue}
\textbf{Step 7 --- ALGEBRA.}

\[
4(2)-\frac4{2^2}=8-\frac4{4}=8-1=7.
\]

\endgroup
\textbf{Result.} The derivative is \(4x-4/x^2\) for \(x\ne0\), and the slope at \(2\)
is the scalar number \(7\).

\textbf{Step 8 --- CHECK: verify the slope with a difference quotient.} Declare \(h\)
as a nonzero scalar increment tending to zero, small enough that \(2+h\ne0\).
Using the equivalent formula, \(f(2)=2(2)^2-3+4/2=8-3+2=7\).

\begingroup\color{TrickPurple}
\textbf{Step 9 --- TRICK: expand the square at the shifted input.} See
T-003.

\[
(2+h)^2=2^2+2(2)h+h^2=4+4h+h^2.
\]

\endgroup
\begingroup\color{AlgebraBlue}
\textbf{Step 10 --- ALGEBRA: form and simplify the quotient.}

\[
\begin{aligned}
f(2+h)&=2(4+4h+h^2)-3+\frac4{2+h}\\
&=8+8h+2h^2-3+\frac4{2+h}\\
&=5+8h+2h^2+\frac4{2+h},\\
f(2+h)-f(2)&=5+8h+2h^2+\frac4{2+h}-7\\
&=8h+2h^2+\frac4{2+h}-2\\
&=8h+2h^2+\frac{4-2(2+h)}{2+h}\\
&=8h+2h^2+\frac{4-4-2h}{2+h}\\
&=8h+2h^2-\frac{2h}{2+h},\\
\frac{f(2+h)-f(2)}h
&=\frac{h(8+2h-2/(2+h))}{h}\\
&=8+2h-\frac2{2+h},\qquad h\ne0.
\end{aligned}
\]

\endgroup
\textbf{Step 11 --- CHECK: take the limit.}

\[
\lim_{h\to0}\frac{f(2+h)-f(2)}h
=8+0-\frac22=8-1=7.
\]

The last arithmetic is \textbf{ALGEBRA}. This independently verifies the requested
slope at \(2\), while the differentiation rules supply the general formula.

\textbf{Recognition cue.} A rational-looking expression may become a sum of powers.
Inspect the algebra before selecting a differentiation rule.

\section{What to carry into the next chapter}

\begin{itemize}
\item A changing inner expression needs its own derivative; see the chain-rule pilot.
\item A product of changing factors requires a product rule unless you first expand it validly.
\item Fixed parameters remain fixed only when the setup says so.
\item A simplified formula does not automatically enlarge the original domain.
\item A check at one input verifies that input, not every value of the derivative.

\end{itemize}
The next chapter will develop product, quotient, and chain-rule choices using
these same declarations and visible algebra steps.

\chapter{Pilot problems and worked solutions}

These examples use notation standard version 0.2 and test the solution granularity. Textual labels indicate the future color styling. Question and solution sections are separate here; typesetting will place solutions on separate pages.

\section{Problems}

\subsection{D2-001 --- Differentiate a composition}

\textbf{Objects and dependencies.} \(x\in\mathbb R\) is an independent scalar variable. The scalar-valued function \(f:\mathbb R\to\mathbb R\) is

\[
f(x)=(3x^2+1)^4.
\]

There are no other variables or parameters. Find the scalar-valued derivative

\[
\frac{d}{dx}\left(f(x)\right).
\]

\subsection{I2-001 --- Evaluate a definite integral}

\textbf{Objects and dependencies.} \(x\) is a scalar integration variable in \([0,1]\). Define the scalar-valued function \(q(x)=2x(x^2+1)^3\). All quantities are real. Integrate with respect to \(x\), with lower limit \(x=0\) and upper limit \(x=1\). Find the scalar number

\[
\int_{0}^{1} q(x)\,dx
=\int_{0}^{1} 2x(x^2+1)^3\,dx.
\]

\section{Hint ladders}

\subsection{D2-001}

\begin{enumerate}
\item Which expression is raised to the fourth power?
\item Define the inner function \(g(x)=3x^2+1\), and the outer function \(h(u)=u^4\). Then \(f(x)=h(g(x))\).
\item Differentiate \(h(u)\) with respect to \(u\), evaluate at \(u=g(x)\), and multiply by the derivative of \(g(x)\) with respect to \(x\).

\end{enumerate}
\subsection{I2-001}

\begin{enumerate}
\item Compare the factor \(2x\) with the derivative of the expression \(x^2+1\).
\item Try the new scalar coordinate \(u=g(x)=x^2+1\).
\item The original endpoints \(x=0\) and \(x=1\) become \(u=1\) and \(u=2\). The transformed integrand is \(u^3\).

\end{enumerate}
\clearpage
\section{Complete solutions}

\subsection{D2-001}

\textbf{Step 1 --- SETUP.} The independent input \(x\) is scalar. The output \(f(x)=(3x^2+1)^4\) is scalar and is defined for all real \(x\). Introduce scalar-valued functions

\[
g(x)=3x^2+1,\qquad h(u)=u^4,\qquad f(x)=h(g(x)).
\]

Here \(u\) is an independent scalar input when defining and differentiating \(h(u)\). In the composition, set \(u=g(x)\), so the inner value depends on \(x\).

\textbf{Step 2 --- CHOICE.} A fourth power surrounds a changing inner expression. The chain rule accounts for both the outer power and the inner rate of change.

\textbf{Step 3 --- CALCULUS: differentiate the outer function.} The power rule gives

\[
\frac{d}{du}\left(h(u)\right)
=\frac{d}{du}\left(u^4\right)
=4u^3.
\]

\textbf{Step 4 --- CALCULUS: differentiate the inner function.} Use the sum rule, constant-multiple rule, power rule, and derivative of a constant:

\[
\begin{aligned}
\frac{d}{dx}\left(g(x)\right)
&=\frac{d}{dx}\left(3x^2+1\right)\\
&=3\frac{d}{dx}\left(x^2\right)+\frac{d}{dx}\left(1\right)\\
&=3(2x)+0.
\end{aligned}
\]

\begingroup\color{AlgebraBlue}
\textbf{Step 5 --- ALGEBRA.} Multiply the numerical factors and remove the zero:

\[
3(2x)+0=(3\cdot2)x+0=6x+0=6x.
\]

\endgroup
\textbf{Step 6 --- CALCULUS: assemble the chain rule.}

\[
\begin{aligned}
\frac{d}{dx}\left(f(x)\right)
&=\left.\frac{d}{du}\left(h(u)\right)\right|_{u=g(x)}
\frac{d}{dx}\left(g(x)\right)\\
&=\left.\left(4u^3\right)\right|_{u=3x^2+1}(6x)\\
&=4(3x^2+1)^3(6x).
\end{aligned}
\]

\begingroup\color{AlgebraBlue}
\textbf{Step 7 --- ALGEBRA.} All factors are scalar, so they may be reordered:

\[
4(3x^2+1)^3(6x)=(4\cdot6)x(3x^2+1)^3=24x(3x^2+1)^3.
\]

\endgroup
\textbf{Result.}

\[
\boxed{\frac{d}{dx}\left(f(x)\right)=24x(3x^2+1)^3.}
\]

\textbf{Step 8 --- CHECK.} At \(x=0\), the formula gives \(24(0)(1)^3=0\). Independently, the derivative definition at zero is

\[
\lim_{\varepsilon\to0}\frac{f(\varepsilon)-f(0)}{\varepsilon}.
\]

Here \(\varepsilon\) is a nonzero scalar increment tending to zero. For the check, the binomial identity
\((a+b)^4=a^4+4a^3b+6a^2b^2+4ab^3+b^4\), with \(a=1\), \(b=3\varepsilon^2\), gives the following \textbf{IDENTITY and ALGEBRA} steps:

\[
\begin{aligned}
f(\varepsilon)&=1+4(3\varepsilon^2)+6(3\varepsilon^2)^2+4(3\varepsilon^2)^3+(3\varepsilon^2)^4\\
&=1+12\varepsilon^2+54\varepsilon^4+108\varepsilon^6+81\varepsilon^8,\\
f(0)&=1,\\
\frac{f(\varepsilon)-f(0)}{\varepsilon}
&=\frac{12\varepsilon^2+54\varepsilon^4+108\varepsilon^6+81\varepsilon^8}{\varepsilon}\\
&=12\varepsilon+54\varepsilon^3+108\varepsilon^5+81\varepsilon^7.
\end{aligned}
\]

Each term tends to zero, so the limit is zero. This checks the answer at one input; it is not a separate proof of the derivative formula at every input.

\textbf{Recognition cue.} Find the outside operation, then identify every input it receives. The inner derivative is required even when the inner expression looks simple.

\subsection{I2-001}

\textbf{Step 1 --- SETUP.} The original scalar variable is \(x\in[0,1]\), and the scalar integrand is \(q(x)=2x(x^2+1)^3\). Introduce

\[
g(x)=x^2+1,\qquad u=g(x),\qquad h(u)=u^3.
\]

Both \(g\) and \(h\) are scalar-valued functions. The new coordinate \(u\) depends on \(x\) until we rewrite the integral in terms of \(u\). Polynomials have the continuity and differentiability needed for the substitution theorem.

\textbf{Step 2 --- CHOICE.} The integrand contains a power of \(x^2+1\) and the factor \(2x\), which is the derivative of \(x^2+1\). This is the reverse of the chain-rule pattern.

\textbf{Step 3 --- CALCULUS.}

\[
\begin{aligned}
\frac{d}{dx}\left(g(x)\right)
&=\frac{d}{dx}\left(x^2+1\right)\\
&=\frac{d}{dx}\left(x^2\right)+\frac{d}{dx}\left(1\right)\\
&=2x+0=2x.
\end{aligned}
\]

The last removal of \(+0\) is an \textbf{ALGEBRA} simplification.

\begingroup\color{AlgebraBlue}
\textbf{Step 4 --- ALGEBRA: match the integrand.} Since the factors are scalars,

\[
2x(x^2+1)^3=(x^2+1)^3(2x)
=h(g(x))\frac{d}{dx}\left(g(x)\right).
\]

\endgroup
\begingroup\color{AlgebraBlue}
\textbf{Step 5 --- ALGEBRA: transform both bounds.}

\[
\begin{aligned}
x=0&\quad\Longrightarrow\quad u=g(0)=0^2+1=0+1=1,\\
x=1&\quad\Longrightarrow\quad u=g(1)=1^2+1=1+1=2.
\end{aligned}
\]

\endgroup
\textbf{Step 6 --- CALCULUS: apply the substitution theorem.}

\[
\begin{aligned}
\int_{0}^{1} 2x(x^2+1)^3\,dx
&=\int_{0}^{1}
\left(h(g(x))\frac{d}{dx}\left(g(x)\right)\right)\,dx\\
&=\int_{g(0)}^{g(1)} h(u)\,du\\
&=\int_{1}^{2} u^3\,du.
\end{aligned}
\]

The whole integral changes coordinates, including the integration variable and both bounds. The derivative factor is accounted for by the theorem.

\textbf{Step 7 --- CALCULUS: find an antiderivative.} The power rule for integration is

\[
\int u^p\,du=\frac{u^{p+1}}{p+1}+C,\qquad p\ne-1,
\]

on an interval where the powers and this rule are defined. Here \(p=3\) is a fixed scalar exponent and \(u\in[1,2]\). Choose one antiderivative \(H(u)\):

\[
H(u)=\frac{u^{3+1}}{3+1}=\frac{u^4}{4}.
\]

The exponent and denominator simplifications are \textbf{ALGEBRA}. A definite integral needs one antiderivative; any added constant cancels in the endpoint difference.

\textbf{Step 8 --- CALCULUS: evaluate at the endpoints.}

\[
\int_{1}^{2} u^3\,du=H(2)-H(1)=\frac{2^4}{4}-\frac{1^4}{4}.
\]

\begingroup\color{AlgebraBlue}
\textbf{Step 9 --- ALGEBRA.}

\[
\begin{aligned}
2^4&=2\cdot2\cdot2\cdot2=16,\\
1^4&=1\cdot1\cdot1\cdot1=1,\\
\frac{2^4}{4}-\frac{1^4}{4}
&=\frac{16}{4}-\frac{1}{4}
=\frac{16-1}{4}
=\frac{15}{4}.
\end{aligned}
\]

\endgroup
\textbf{Result.}

\[
\boxed{\int_{0}^{1} 2x(x^2+1)^3\,dx=\frac{15}{4}.}
\]

\textbf{Step 10 --- CHECK: differentiate the original-variable antiderivative.} Define the scalar-valued function \(Q(x)=(x^2+1)^4/4\). The constant-multiple and chain rules give

\[
\frac{d}{dx}\left(Q(x)\right)
=\frac14\,4(x^2+1)^3(2x).
\]

\textbf{ALGEBRA:} \((1/4)\cdot4=1\), so this becomes \(2x(x^2+1)^3=q(x)\), the original integrand. Also,

\[
Q(1)-Q(0)=\frac{(1^2+1)^4}{4}-\frac{(0^2+1)^4}{4}
=\frac{2^4}{4}-\frac{1^4}{4}=\frac{15}{4}.
\]

\textbf{Recognition cue.} Look for an inner expression whose derivative appears as a multiplying factor. A near match may need a constant adjustment, which must be written explicitly.

\section{Tutor entry point}

Example request: “Help me with I2-001, Step 6. I understand the new bounds but not what happened to the factor \(2x\). Use the notation standard and give me one step at a time.”

\section{Review questions before drafting full chapters}

\begin{itemize}
\item Are the round grouping parentheses and nearby “held fixed” statements easy to follow?
\item Do the pilot algebra steps give enough detail, or should common arithmetic move into expandable side explanations in a digital edition?
\item Does the outer-function/inner-function setup make substitution and the chain rule easier to follow?

\end{itemize}
Practice variants and their complete solutions will be added when these pilots become chapter material.

\appendix
\chapter{Tricks and identities}
\begingroup\color{TrickPurple}

Alpha reference appendix. These are supporting identities, algebraic rewrites,
and approximation tools to consult while working through problems. Main calculus
methods such as substitution, integration by parts, and differentiation rules
belong in the main chapters.

\textbf{Purple means a supporting trick.} The label TRICK remains visible in grayscale.
The main chapters use dark blue for algebra, black for calculus rules, and purple
when invoking a supporting tool from this appendix.

\section{Find a useful tool}

\begin{longtable}{@{}>{\raggedright\arraybackslash}p{0.300\textwidth}>{\raggedright\arraybackslash}p{0.300\textwidth}>{\raggedright\arraybackslash}p{0.300\textwidth}@{}}
Tool & When to consider it & What it gives \\[5pt]
\hline\endhead
T-001 Euler identities & Trigonometric products or oscillations & Exact conversions between complex exponentials and real trigonometric functions \\[5pt]
T-002 Taylor expansions & A local approximation near a known input & A polynomial plus an error term; an infinite identity only when convergence is justified \\[5pt]
T-003 Binomial expansion & Powers of a sum & A finite exact expansion for nonnegative integer powers \\[5pt]
T-004 Completing the square & A quadratic is hard to interpret or manipulate & An exact shifted-square form \\[5pt]
T-005 Rationalization & A difference of roots causes cancellation & An exact equivalent expression on a stated domain \\[5pt]
T-006 Logarithm identities & Products, quotients, and powers inside a logarithm & Exact sums and differences under real positivity assumptions \\[5pt]
\end{longtable}

\section{T-001 Euler identities}

\textbf{Objects and conditions.} \(\theta\in\mathbb R\) is an independent real scalar
angle, measured in radians. \(\mathrm i\) is the fixed imaginary scalar with
\(\mathrm i^2=-1\). A complex scalar \(z=a+\mathrm i b\) has real scalar parts \(a,b\);
\(\operatorname{Re}(z)=a\) and \(\operatorname{Im}(z)=b\).
The complex exponential \(\exp(z)\) is a scalar-valued function of a complex scalar
input. Sine and cosine of real angles are real scalars.
The usual notation \(e^z\) means the same function as \(\exp(z)\), where
\(e=\exp(1)\) is the fixed real Euler constant. We use \(\exp(z)\) here to keep the
input visible.

\textbf{TRICK --- exact identities.}

\[
\exp(\mathrm i\theta)=\cos(\theta)+\mathrm i\sin(\theta),
\qquad
\exp(-\mathrm i\theta)=\cos(\theta)-\mathrm i\sin(\theta).
\]

Adding and subtracting these equations gives

\[
\cos(\theta)=\frac{\exp(\mathrm i\theta)+\exp(-\mathrm i\theta)}{2},
\qquad
\sin(\theta)=\frac{\exp(\mathrm i\theta)-\exp(-\mathrm i\theta)}{2\mathrm i}.
\]

At \(\theta=\pi\), where \(\pi\) is the fixed real circle constant,
\(\cos(\pi)=-1\) and \(\sin(\pi)=0\), so \(\exp(\mathrm i\pi)+1=0\).
These formulas are recorded in \href{https://dlmf.nist.gov/4.2}{NIST DLMF §4.2}.

\subsection{Worked use: derive double-angle identities}

\textbf{Step 1 --- SETUP.} The same real scalar \(\theta\) supplies every angle below.
All expressions containing \(\mathrm i\) are complex scalars.

\begingroup\color{TrickPurple}
\textbf{Step 2 --- TRICK: use the exponential addition identity.} The complex
exponential satisfies \(\exp(z+w)=\exp(z)\exp(w)\) for complex scalar inputs \(z,w\).
With \(z=w=\mathrm i\theta\),

\[
\exp(2\mathrm i\theta)
=\exp(\mathrm i\theta)\exp(\mathrm i\theta)
=\left(\cos(\theta)+\mathrm i\sin(\theta)\right)^2.
\]

\endgroup
\begingroup\color{TrickPurple}
\textbf{Step 3 --- ALGEBRA: expand and use the imaginary unit's square.}

\[
\begin{aligned}
\left(\cos(\theta)+\mathrm i\sin(\theta)\right)^2
&=(\cos(\theta))^2
+2\mathrm i\cos(\theta)\sin(\theta)
+\mathrm i^2(\sin(\theta))^2\\
&=(\cos(\theta))^2-(\sin(\theta))^2
+\mathrm i\left(2\cos(\theta)\sin(\theta)\right).
\end{aligned}
\]

\endgroup
\begingroup\color{TrickPurple}
\textbf{Step 4 --- TRICK: compare the real and imaginary parts.} Euler's formula also
gives \(\exp(2\mathrm i\theta)=\cos(2\theta)+\mathrm i\sin(2\theta)\).
Equal complex numbers have equal real parts and equal imaginary parts:

\[
\cos(2\theta)=(\cos(\theta))^2-(\sin(\theta))^2,
\qquad
\sin(2\theta)=2\cos(\theta)\sin(\theta).
\]

\endgroup
\textbf{Step 5 --- CHECK.} At \(\theta=0\), the cosine formula gives \(1=1-0\) and the
sine formula gives \(0=2(1)(0)\). This checks one angle. The preceding exact
identities establish the general result.

\textbf{When it helps.} Converting trigonometric expressions into exponentials can
make angle combinations easier to manipulate. Track real and imaginary parts
explicitly when returning to a real-valued problem.

\section{T-002 Taylor expansions}

\textbf{Objects and conditions.} \(x\) is an independent real scalar and \(a\) is a fixed
real expansion center. \(f(x)\) is a real scalar-valued function. \(N\) is a fixed
nonnegative integer polynomial degree. \(k\) is an integer summation index;
\(t\) is an auxiliary real scalar input used when writing derivatives before
evaluation. Both \(f(t)\) and \(f(x)\) are values of the same function at their
displayed inputs. The factorial is \(k!=1\cdot2\cdots k\) for \(k\ge1\), with \(0!=1\).

\textbf{TRICK --- Taylor polynomial.} Define the fixed scalar coefficients by

\[
c_0=f(a),\qquad
c_k=\frac1{k!}\left.\frac{d^k}{dt^k}\left(f(t)\right)\right|_{t=a},
\qquad 1\leq k\leq N.
\]

Then the scalar polynomial is

\[
P_N(x;a)=\sum_{k=0}^{N}c_k(x-a)^k.
\]

For \(N=3\), the full expression is

\[
\begin{aligned}
P_3(x;a)={}&f(a)
+\left.\frac{d}{dt}\left(f(t)\right)\right|_{t=a}(x-a)\\
&+\frac1{2!}\left.\frac{d^2}{dt^2}\left(f(t)\right)\right|_{t=a}(x-a)^2\\
&+\frac1{3!}\left.\frac{d^3}{dt^3}\left(f(t)\right)\right|_{t=a}(x-a)^3.
\end{aligned}
\]

A Maclaurin expansion is the special case \(a=0\). The superscript on a derivative
counts repeated differentiation; it does not square or cube the derivative.

\subsection{The error is part of the tool}

If \(f\) has continuous derivatives through order \(N+1\) on an interval containing
\(a\) and \(x\), Taylor's theorem gives a scalar remainder \(R_N(x;a)\):

\[
f(x)=P_N(x;a)+R_N(x;a),
\]

\[
R_N(x;a)=\frac1{(N+1)!}
\left.\frac{d^{N+1}}{dt^{N+1}}\left(f(t)\right)\right|_{t=\xi}
(x-a)^{N+1},
\]

for some real scalar \(\xi\) between \(a\) and \(x\) when they differ. If a fixed
nonnegative real scalar \(M\) bounds the absolute value of that derivative
throughout the intervening interval, then

\[
|R_N(x;a)|\leq\frac{M|x-a|^{N+1}}{(N+1)!}.
\]

An infinite Taylor series equals \(f(x)\) only where the remainder tends to zero
as \(N\) increases. Smoothness alone does not guarantee that equality. Keep a
finite approximation marked \(\approx\) unless the remainder is included.
See \href{https://openstax.org/books/calculus-volume-2/pages/6-3-taylor-and-maclaurin-series}{OpenStax, Calculus Volume 2, §6.3}.

\subsection{Common infinite expansions and their valid inputs}

Here \(x\) is an independent real scalar, and \(k\) is an integer summation index.
Each displayed function has a real scalar output. For every real \(x\),

\[
\begin{aligned}
\exp(x)&=\sum_{k=0}^{\infty}\frac{x^k}{k!}
=1+x+\frac{x^2}{2!}+\frac{x^3}{3!}+\cdots,\\
\sin(x)&=\sum_{k=0}^{\infty}\frac{(-1)^k x^{2k+1}}{(2k+1)!}
=x-\frac{x^3}{3!}+\frac{x^5}{5!}-\cdots,\\
\cos(x)&=\sum_{k=0}^{\infty}\frac{(-1)^k x^{2k}}{(2k)!}
=1-\frac{x^2}{2!}+\frac{x^4}{4!}-\cdots.
\end{aligned}
\]

The sine and cosine inputs are angles in radians. The summation formulas
specify every term; the trailing dots only illustrate their pattern. These
series converge to their displayed functions at every real input. Keeping
only finitely many terms requires an approximation sign and an error bound.

The geometric expansion has a narrower range of valid inputs:

\[
\frac{1}{1-x}=\sum_{k=0}^{\infty}x^k
=1+x+x^2+x^3+\cdots,\qquad |x|<1.
\]

Although the function \(1/(1-x)\) exists at \(x=2\), this series does not converge
there. Check the series' convergence range before replacing a function by a
series. See \href{https://openstax.org/books/calculus-volume-2/pages/6-4-working-with-taylor-series}{OpenStax, Calculus Volume 2, §6.4}.

\subsection{Worked use: approximate a sine with a stated error}

\textbf{Step 1 --- SETUP.} \(f(x)=\sin(x)\) is a real scalar function, with real scalar
angle \(x\) in radians. Use center \(a=0\) and degree \(N=4\). The target input is
\(x=1/10\) radians. There are no varying parameters.

\textbf{Step 2 --- CALCULUS: list the derivatives used by the tool.}

\[
\begin{aligned}
f(t)&=\sin(t),\\
\frac{d}{dt}\left(f(t)\right)&=\cos(t),\\
\frac{d^2}{dt^2}\left(f(t)\right)&=-\sin(t),\\
\frac{d^3}{dt^3}\left(f(t)\right)&=-\cos(t),\\
\frac{d^4}{dt^4}\left(f(t)\right)&=\sin(t),\\
\frac{d^5}{dt^5}\left(f(t)\right)&=\cos(t).
\end{aligned}
\]

At \(t=0\), the value and first four derivatives are respectively \(0,1,0,-1,0\).
These trigonometric derivative rules are developed in the main calculus material;
here their values supply the expansion coefficients.

\begingroup\color{TrickPurple}
\textbf{Step 3 --- ALGEBRA: build the polynomial and remove zero terms.}

\[
\begin{aligned}
P_4(x;0)
&=0+1(x-0)+\frac0{2!}(x-0)^2
+\frac{-1}{3!}(x-0)^3+\frac0{4!}(x-0)^4\\
&=x-\frac{x^3}{1\cdot2\cdot3}
=x-\frac{x^3}{6}.
\end{aligned}
\]

Although this polynomial's highest nonzero power is three, it is also the
fourth-order Taylor polynomial because its fourth-order coefficient is zero.

\endgroup
\begingroup\color{TrickPurple}
\textbf{Step 4 --- APPROXIMATION: retain the fifth-derivative error bound.} Since
\(|\cos(t)|\leq1\) for real \(t\), choose the fixed scalar bound \(M=1\):

\[
\sin(x)=x-\frac{x^3}{6}+R_4(x;0),
\qquad |R_4(x;0)|\leq\frac{|x|^5}{5!}.
\]

\endgroup
\begingroup\color{TrickPurple}
\textbf{Step 5 --- ALGEBRA: evaluate the polynomial exactly.}

\[
\begin{aligned}
P_4(1/10;0)
&=\frac1{10}-\frac{(1/10)^3}{6}\\
&=\frac1{10}-\frac1{1000\cdot6}\\
&=\frac{600}{6000}-\frac1{6000}
=\frac{599}{6000}.
\end{aligned}
\]

\endgroup
\begingroup\color{TrickPurple}
\textbf{Step 6 --- APPROXIMATION: report the result with its bound.}

\[
\sin(1/10)\approx\frac{599}{6000},
\qquad
\left|\sin(1/10)-\frac{599}{6000}\right|
\leq\frac1{10^5\cdot120}
=\frac1{12\,000\,000}.
\]

This rational bound applies to the exact rational approximation shown. If you
round it to a decimal, include the additional rounding error in the reported bound.

\endgroup
\textbf{Step 7 --- CHECK.} The correction \(x^3/6\) is positive at \(x=1/10\), so the
polynomial estimate is slightly smaller than \(x\). Also, both the polynomial and
the actual sine are zero at the expansion center. These plausibility checks
supplement the theorem's error bound.

\textbf{When it helps.} Expand around an input where function values and derivatives
are simple. A polynomial without an error discussion is an incomplete numerical
approximation.

\section{T-003 Binomial expansion}

\textbf{Objects and conditions.} \(r,s\) are real or complex scalar expressions.
\(n\) is a fixed nonnegative integer, and \(k\) is an integer index from zero to \(n\).
The binomial coefficient is the scalar integer

\[
\binom nk=\frac{n!}{k!(n-k)!}.
\]

\textbf{TRICK --- exact finite identity.}

\[
(r+s)^n=\sum_{k=0}^n\binom nk r^{n-k}s^k.
\]

For a zero power, the corresponding factor is \(1\); the endpoint terms are
\(r^n\) and \(s^n\). This is a finite polynomial identity. A fractional exponent
requires a different, generally infinite expansion with convergence conditions.

\subsection{Worked use: expand a shifted cubic}

\textbf{Step 1 --- SETUP.} \(x\) is an independent real scalar. \(h\) is a real scalar
increment. Here take \(r=x\), \(s=h\), and \(n=3\).

\begingroup\color{TrickPurple}
\textbf{Step 2 --- ALGEBRA: calculate the coefficients.}

\[
\begin{aligned}
\binom30&=\frac{3!}{0!3!}=1,\\
\binom31&=\frac{3!}{1!2!}=\frac{6}{1\cdot2}=3,\\
\binom32&=\frac{3!}{2!1!}=\frac{6}{2\cdot1}=3,\\
\binom33&=\frac{3!}{3!0!}=1.
\end{aligned}
\]

\endgroup
\begingroup\color{TrickPurple}
\textbf{Step 3 --- TRICK: write every term.}

\[
\begin{aligned}
(x+h)^3
&=1(x^3)+3(x^2h)+3(xh^2)+1(h^3)\\
&=x^3+3x^2h+3xh^2+h^3.
\end{aligned}
\]

\endgroup
\textbf{Step 4 --- CHECK.} At \(h=0\) the expression reduces to \(x^3\); at \(x=0\) it
reduces to \(h^3\). These checks do not replace the binomial identity but can expose
a missing endpoint term. For a square, the same tool gives
\((x+h)^2=x^2+2xh+h^2\).

\textbf{When it helps.} Shifted inputs occur in derivative definitions and series
calculations. The coefficients prevent expansion errors without hiding terms.

\section{T-004 Completing the square}

\textbf{Objects and conditions.} \(x\) is an independent real scalar; \(a,b,c\) are fixed
real scalar coefficients, with \(a\ne0\). The scalar expression is \(ax^2+bx+c\).

\textbf{TRICK --- exact identity.}

\[
ax^2+bx+c=a\left(x+\frac{b}{2a}\right)^2+c-\frac{b^2}{4a}.
\]

\subsection{Worked use: expose a quadratic's shape}

\textbf{Step 1 --- SETUP.} Let the real scalar function be \(q(x)=2x^2-8x+11\) for all
real \(x\). The coefficients \(2,-8,11\) are fixed.

\begingroup\color{TrickPurple}
\textbf{Step 2 --- ALGEBRA: factor the quadratic coefficient.}

\[
q(x)=2(x^2-4x)+11.
\]

\endgroup
\begingroup\color{TrickPurple}
\textbf{Step 3 --- TRICK: add and subtract the missing square.} Since
\((x-2)^2=x^2-4x+4\),

\[
x^2-4x=(x-2)^2-4.
\]

\endgroup
\begingroup\color{TrickPurple}
\textbf{Step 4 --- ALGEBRA: substitute and simplify.}

\[
\begin{aligned}
q(x)&=2\left((x-2)^2-4\right)+11\\
&=2(x-2)^2-8+11\\
&=2(x-2)^2+3.
\end{aligned}
\]

\endgroup
\textbf{Step 5 --- CHECK: expand back.}

\[
2(x-2)^2+3=2(x^2-4x+4)+3
=2x^2-8x+8+3=2x^2-8x+11.
\]

The square is nonnegative, so \(q(x)\geq3\), with equality at \(x=2\). The exact
rewrite exposes that feature without requiring an optimization calculation.

\textbf{When it helps.} A shifted-square denominator or root may reveal the right
identity, domain, or later integration method. The rewrite itself does not choose
the calculus method for you.

\section{T-005 Rationalization}

\textbf{Objects and conditions.} If \(r,s\) are real scalar expressions, the exact
difference-of-squares identity is

\[
(r-s)(r+s)=r^2-s^2.
\]

The factor \(r+s\) is often called the conjugate of \(r-s\) in this real-root
algebra context. Multiplying a fraction by \((r+s)/(r+s)\) preserves it only when
\(r+s\ne0\).

\subsection{Worked use: remove a root difference}

\textbf{Step 1 --- SETUP.} \(x\) is an independent real scalar with \(x>-1\) and \(x\ne0\).
Define the real scalar function

\[
f(x)=\frac{\sqrt{1+x}-1}{x}.
\]

The square root is the nonnegative real root. On this domain,
\(\sqrt{1+x}+1>0\), so the conjugate multiplier is not zero.

\begingroup\color{TrickPurple}
\textbf{Step 2 --- TRICK: multiply by one in a useful form.}

\[
f(x)=\frac{\sqrt{1+x}-1}{x}
\frac{\sqrt{1+x}+1}{\sqrt{1+x}+1}.
\]

\endgroup
\begingroup\color{TrickPurple}
\textbf{Step 3 --- ALGEBRA: multiply the numerators and cancel on the allowed domain.}

\[
\begin{aligned}
f(x)&=\frac{(\sqrt{1+x}-1)(\sqrt{1+x}+1)}{x(\sqrt{1+x}+1)}\\
&=\frac{(\sqrt{1+x})^2-1^2}{x(\sqrt{1+x}+1)}\\
&=\frac{1+x-1}{x(\sqrt{1+x}+1)}\\
&=\frac{x}{x(\sqrt{1+x}+1)}\\
&=\frac1{\sqrt{1+x}+1},\qquad x>-1,\quad x\ne0.
\end{aligned}
\]

\endgroup
\textbf{Step 4 --- CHECK.} At the allowed input \(x=3\), the original expression gives
\((\sqrt4-1)/3=(2-1)/3=1/3\); the rewritten expression gives \(1/(\sqrt4+1)=1/3\).
The original function remains undefined at \(x=0\). The rewritten expression has
a value there and supplies a continuous extension, not a new value of the original
function. See D1-004 for the same domain distinction.

\textbf{When it helps.} A root difference may be hard to simplify or numerically
evaluate near cancellation. Rationalization preserves exact equality where its
divisions are allowed. It can prepare an expression for a limit calculation.

\section{T-006 Logarithm identities}

\textbf{Objects and conditions.} \(r,s\) are positive real scalar expressions and \(p\)
is a fixed real scalar exponent. The natural logarithm \(\ln(r)\) is real. On
these real domains,

\[
\ln(rs)=\ln(r)+\ln(s),\qquad
\ln\left(\frac r s\right)=\ln(r)-\ln(s),\qquad
\ln(r^p)=p\ln(r).
\]

These identities need positivity in this real formulation. Complex logarithms
require branch choices and do not permit these unrestricted rewrites. See the
exponential and logarithm definitions in \href{https://dlmf.nist.gov/4.2}{NIST DLMF §4.2}.

\subsection{Worked use: expose a product and power inside a logarithm}

\textbf{Step 1 --- SETUP.} \(x>0\) is an independent real scalar. The real scalar
function is

\[
q(x)=\ln\left(\frac{(x^2+1)^3}{x^2}\right).
\]

On this domain, \(x^2+1>0\) and \(x^2>0\), so both logarithms introduced below
have positive arguments.

\begingroup\color{TrickPurple}
\textbf{Step 2 --- TRICK: separate the quotient.}

\[
q(x)=\ln\left((x^2+1)^3\right)-\ln(x^2).
\]

\endgroup
\begingroup\color{TrickPurple}
\textbf{Step 3 --- TRICK: move the fixed powers outside.}

\[
q(x)=3\ln(x^2+1)-2\ln(x),\qquad x>0.
\]

\endgroup
\textbf{Step 4 --- CHECK.} At \(x=1\), the original expression is
\(\ln((1+1)^3/1)=\ln(8)\). The new expression is \(3\ln(2)-2\ln(1)=3\ln(2)\),
which equals \(\ln(8)\) by the power identity, with \(\ln(1)=0\).

\textbf{When it helps.} Revealing sums and fixed multipliers can organize a later
derivative calculation. This algebraic rewrite does not eliminate the dependence
of \(x^2+1\) on \(x\).

For a version declared on \(x\ne0\), including negative inputs, the last expression
would instead be \(3\ln(x^2+1)-2\ln(|x|)\). The absolute value keeps the real
logarithm's input positive; the expression \(\ln(x)\) must not be carried into a
negative-input domain.

\section{Use a trick deliberately}

Identify the structure it matches, check its domain, show the exact replacement,
and then return to the calculus problem. Keep an error term or error bound when
using an approximation. If a trick makes the expression longer or obscures its
dependencies, keep the original form and choose another approach.
\endgroup

\chapter{Literature review and notation decisions}

Reviewed October 5, 2026. Scope: selected introductory calculus textbook sections, a university matrix-calculus course, a matrix-calculus reference, a tensor survey, and continuum-mechanics notes. This is a targeted notation review, not an exhaustive survey or a study demonstrating that one notation improves learning for everyone.

\section{Recommendation}

Use established Leibniz operator notation with round grouping parentheses and complete function arguments. Pair it with local type declarations, dependency statements, and explicit held-fixed instructions. The symbols should transfer to other books; the extra explanation should reduce reliance on memory.

The recommendation is an editorial judgment tailored to this reader. The references establish usage and mathematical distinctions, not a uniquely “best” personal notation.

\section{Findings and decisions}

\subsection{R1 --- Ordinary derivatives}

OpenStax explicitly lists the operator form \(\frac{d}{dx}\left(f(x)\right)\) alongside prime and quotient-style notation. It also discusses higher derivatives. \textbf{Decision:} adopt that operator form; retain \(d^2/dx^2\) with an initial nested expansion. Prime and dot notation belong in a translation guide, not worked calculations. \href{https://openstax.org/books/calculus-volume-1/pages/3-2-the-derivative-as-a-function}{OpenStax, Calculus Volume 1, §3.2}.

\subsection{R2 --- Partial derivatives}

OpenStax defines partial derivatives by changing one coordinate while the other stays fixed, and uses both fraction and subscript notation. \textbf{Decision:} retain the standard partial symbol, show all function arguments, and state the held-fixed inputs in nearby text. Our insistence on repeating the instruction is a teaching convention. \href{https://openstax.org/books/calculus-volume-3/pages/4-3-partial-derivatives}{OpenStax, Calculus Volume 3, §4.3}.

\subsection{R3 --- Dependencies and the chain rule}

OpenStax's multivariable chain rule distinguishes independent, intermediate, and dependent variables and uses dependency trees. \textbf{Decision:} show \(z(t)=f(x(t),y(t))\) explicitly and show where each partial derivative is evaluated. A partial derivative in independent coordinates must not be confused with a derivative along a path. \href{https://openstax.org/books/calculus-volume-3/pages/4-5-the-chain-rule}{OpenStax, Calculus Volume 3, §4.5}.

\subsection{R4 --- Substitution and integrals}

OpenStax uses ordinary integral bounds and differentials, supplies a definite-substitution theorem, and demonstrates changing endpoints. \textbf{Decision:} use conventional integral notation with an adjacent variable-and-bounds statement. Explain differential identities before using them. Our earlier blanket exclusion of differential notation was unnecessarily restrictive; fully explained notation can remain standard. \href{https://openstax.org/books/calculus-volume-1/pages/5-5-substitution}{OpenStax, Calculus Volume 1, §5.5}.

\subsection{R5 --- Gradients, Jacobians, and derivative maps}

MIT's course treats derivatives as linear operators. For a map from \(n\) inputs to \(m\) outputs, the Jacobian is \(m\times n\). For a scalar output in Euclidean coordinates, the gradient is a column vector and the derivative is represented by its transpose. \textbf{Decision:} keep these objects distinct, explicitly state shapes, and introduce derivative maps for general matrix-valued functions. We choose \(\mathrm D\) instead of the course's prime notation, with round parentheses for the increment argument. \href{https://ocw.mit.edu/courses/18-s096-matrix-calculus-for-machine-learning-and-beyond-january-iap-2023/pages/lecture-notes-and-readings/}{MIT 18.S096, Lecture Notes and Readings, Lecture 1}.

\subsection{R6 --- Matrix derivatives and constraints}

The Matrix Cookbook defines a scalar function's matrix derivative through derivatives with respect to entries. Its structured-matrix section distinguishes constrained entries from independent entries. \textbf{Decision:} show entry derivatives first, specify the matrix-gradient shape, and name independent parameters for constrained matrices. Do not import an unexplained matrix derivative fraction from another source. \href{https://math.uwaterloo.ca/~hwolkowi/matrixcookbook}{Petersen and Pedersen, The Matrix Cookbook, November 15, 2012, §§2 and 2.8}.

\subsection{R7 --- Tensor typography and index meaning}

Kolda and Bader use bold lowercase vectors, bold uppercase matrices, and bold script higher-order tensors, with ordinary scalar component symbols. They distinguish order from rank and note alternative conventions. \textbf{Decision:} adopt that visual distinction, using the available bold calligraphic font, while explicitly declaring tensor spaces and bases when discussing geometric tensors. Typography alone cannot explain the object. \href{https://www.kolda.net/publication/TensorReview.pdf}{Kolda and Bader, Tensor Decompositions and Applications, SIAM Review 51(3), 2009, §2}.

\subsection{R8 --- Summation shorthand}

Stanford continuum-mechanics notes place explicit sums beside their Einstein-convention equivalents and expand the component equations. \textbf{Decision:} retain explicit sums and limits in our solutions; teach the implicit convention only as a way to read outside references. This is an intentional accessibility choice within ordinary mathematics. \href{https://biomechanics.stanford.edu/me338/me338_n01.pdf}{Stanford ME338A, Continuum Mechanics, Lecture Notes 01, §1.1.1.1}.

\section{Delimiter and symbol policy}

These are editorial decisions informed by the review:

\begin{itemize}
\item Round parentheses group derivative operands. Square brackets display vectors and matrices. Closed intervals retain their conventional square brackets, with endpoint inequalities spelled out when introduced.
\item Tensor components use declared indices. Square brackets are not a universal tensor marker.
\item Ordinary integrands do not receive decorative wrappers. Group sums only where grouping clarifies the expression.
\item Evaluation bars mean substitution at specified inputs. “Held fixed” appears in words, so those two instructions cannot be mistaken for one another.
\item Use \(\arcsin(x)\) for the inverse sine and \(1/\sin(x)\) for its reciprocal. Use \((\sin(x))^2\) to expose a function power.
\item Use \(\det(\mathbf A)\) for a determinant and absolute-value bars for scalar magnitude.
\item A short translation guide will expose outside shorthand without changing the notation used in the book's worked steps.

\end{itemize}
\section{Changes from version 0.1}

Round derivative operands replace square operands; redundant parentheses around the operator are removed. Ordinary integral bounds replace variable-tagged bounds, with explicit endpoint statements retained in prose. Held-fixed prose replaces bulky operator subscripts. Higher-order operator notation is taught through expansion rather than excluded. Differential substitution notation is explained rather than treated as intrinsically suspect. The new standard also addresses inverse, norm, determinant, index, and shape ambiguities.
\end{document}
